\section{Conclusion}
This research evaluates GPT-4o-mini in the role of a mutant selector - a role that has never been studied for LLMs - compared with fixed-seed selection under the same top-10\% budget, on 11,040 PIT mutants \cite{co}
from three Defects4J projects spanning 3 different domains.

Answer to RQ1: "We investigated whether GPT-4o-mini-based mutant selection achieves a higher mutation score than same-budget random selection. Under the pre-registered group-weighted test, the evidence was insufficient to reject H0 (Wilcoxon p=0.059, $r\_rb$=+0.350, N=37 groups); under the overall mutation score, GPT-4o-mini-based selection outperformed random selection by +4.6 points (0.716 vs. 0.670), with p=0.0097 (two-proportion z-test) and p=0.0002 (5,000-draw Monte Carlo)." The overall result is "partial support": the selection signal is real and statistically robust at the mutant level, but it is not uniform across classes - strong in complex parsers (JacksonCore, 9 wins/3 losses), and disappearing in code with near-saturated test suites (Codec, base MS 84.9\%).

Key contributions: this is the first study (i) to use GPT-4o-mini as a decision maker for selecting mutants against random selection under the same budget, and (ii) to show that group-weighted versus mutant-weighted results can reverse the statistical conclusion in selection problems - a specific threat that future studies should proactively report on both sides, continuing the pitfall-systematization spirit of \cite{delgadorerez2026pitfalls}.

Future work: (1) Increase power where needed: expand the number of classes within the same project (running PIT at full scope instead of only on classes affected by the patch) - a direct lever for increasing the number of groups, rather than adding new projects; (2) Test the complexity hypothesis: stratify results by the cyclomatic complexity of each class to directly confirm the "LLM helps most in complex code" pattern suggested by the current data; (3) Increase score resolution: replace the 1-5 scale with pairwise ranking or calibrated kill probability - a direction already signaled by prioritization work \cite{bouafif2025primg} - to reduce ties, the main source of information loss in group-weighted testing.